\subsection{An intrinsic cubic obstruction inside the zero toric fiber}

\begin{statement}{E.31}{Proved}{Two abelian classes in one toric fiber have different liftability}
Use the pro-$3$ version of E.24, with its unit rational tangent
at $O=[1:-1:0]$. Put $K=\mathbb Q(\zeta)$,
$\lambda=1-\zeta$ and $\pi=3+\zeta$, so
$N_{K/\mathbb Q}(\pi)=7$.
In the frame coordinates of (E24.4) take the rational torus point
\[
 t=(7,7,1,1,\pi)\in T_D(\mathbb Q)\subset G(\mathbb Q).
 \tag{E31.1}
\]
Its Kummer class
$\alpha\in H^1(G_{\mathbb Q},T_3G)$ has zero image in the
toric minimum. Let $\mathcal K$ denote the kernel of that toric
map on section classes. Every class in the nonempty coset
$\alpha+3\mathcal K$ fails to lift to the cubic's own class-two
pro-$3$ arithmetic quotient, its metabelian quotient, and its
marked-depth quotient of E.25. The zero class in the same toric
fiber lifts to a full tangential section.

The obstruction is visible in an explicit quotient of the
intrinsic based Galois action over $K$.
Let $H_3$ be the group of order $27$ with coordinates
$(a,b,c)\in\mathbb F_3^3$ and multiplication
\[
 (a,b,c)(a',b',c')
 =\left(a+a',b+b',
 c+c'+\tfrac12(ab'-ba')\right).
 \tag{E31.2}
\]
There is a surjection from the cubic geometric group to $H_3$
whose Galois action, with the tangent at $O$ as origin, is
\[
 g:(a,b,c)\longmapsto
       (a,b,c+\kappa_3(\zeta)(g)b),
       \qquad g\in G_K .
 \tag{E31.3}
\]
Thus (E31.3) gives the arithmetic cross term explicitly, as
well as a finite obstruction to the indicated section classes.
\end{statement}
\begin{proof}
Equation (E24.4) sends (E31.1) to
$(7/N(\pi),7/N(\pi))=(1,1)$.
Kummer functoriality therefore puts $\alpha$ in $\mathcal K$.
This is a genuine geometrically abelian section class by E.24.
It is nonzero modulo three: the valuation at the prime $(\pi)$
in the third $K^\times$ factor of (E24.7) is zero on its
displayed $\lambda$ relation, and takes the value one on $t$.
It extends to an integral homomorphism on that rational-point
presentation, so $t$ is not three times a rational point of $G'$.

Over $K$ use the actual morphism
\[
 h=\frac{x+y}{x+\zeta y}:
 U_3\longrightarrow\mathbb P^1\setminus\{0,1,\infty\},
 \qquad
 1-h=\frac{(\zeta-1)y}{x+\zeta y}.
 \tag{E31.4}
\]
The equations $h=0$ and $h=\infty$ force $z=0$ on the cubic;
$h=1$ forces $y=0$. All these points have been removed from
$U_3$, so the displayed target and morphism are exact.
Write $C_i=[1:-\zeta^i:0]$ and $B_i=[-\zeta^i:0:1]$.
The numerator and denominator tangent lines give
\[
 \operatorname{div}(h)=3O-3C_2,\qquad
 \operatorname{div}(1-h)=B_0+B_1+B_2-3C_2 .
 \tag{E31.5}
\]
The frame at $C_2$ is $\bar\pi$, not $\pi$.
Formula (E24.3) consequently gives the exact character values
\[
 \chi_h(t)=\bar\pi^{-3},\qquad
 \chi_{1-h}(t)=7\bar\pi^{-3}.
 \tag{E31.6}
\]

The based arithmetic map is obtained by applying (E31.4) to
the tangential fiber functor at $O$ and choosing a geometric
path to the standard unit tangent at zero in the target.
Its affine effect can be computed without a path convention.
At $O$ put $s=z/x$; then
$y/x=-(1+s^3)^{1/3}$, and hence
\[
 h=-\frac{s^3}{3(1-\zeta)}+O(s^6)
   =\zeta\lambda^{-3}s^3+O(s^6),
 \qquad 1-h=1+O(s^3).
 \tag{E31.7}
\]
The identity in the second expression is
$\lambda^2=-3\zeta$.
The map of Puiseux fiber functors under (E31.7) has the
constant Kummer factors $(\zeta\lambda^{-3},1)$.
Thus the image of the tangential section at $O$ has, modulo
three, abelian coordinates
$(\kappa_3(\zeta),0)$ relative to the standard target origin.
Combining (E31.6) and (E31.7), the abelian section with class
$\alpha$ has target coordinates
\[
 \bigl(\kappa_3(\zeta\lambda^{-3}\bar\pi^{-3}),
        \kappa_3(7\bar\pi^{-3})\bigr)
   =(\kappa_3(\zeta),\kappa_3(7)).
 \tag{E31.8}
\]
This retains the nonzero affine origin term explicitly.

Take the quotient of the target free pro-$3$ group in which
the two generators and their central commutator have order
three; it is (E31.2). Its standard based Galois action over
$K$ is trivial. Indeed $\mu_3\subset K$, and the normalized
action of E.20 has cyclotomic exponent one modulo three and
a conjugator with central image.
Let $\phi_h:N_3\to H_3$ be the geometric map induced by
(E31.4). It is surjective. The two geometric Kummer characters
$h,1-h$ are independent modulo cubes: the valuation at a
$B_i$ forces the coefficient of $1-h$ to be zero modulo
three; $h$ is not a cube, since such a cube root would have
divisor $O-C_2$, a degree-one map from an elliptic curve to
$\mathbb P^1$. Their independent abelian images generate
$H_3/[H_3,H_3]$, and their commutator generates its central
subgroup, proving surjectivity.

Let $q_g\in H_3$ be the first component of the image of the
tangential section at $O$ in
$H_3\times G_K$. Equivariance of the arithmetic map gives
$\phi_h(gng^{-1})=q_g\phi_h(n)q_g^{-1}$.
Its abelian coordinate is
$(\kappa_3(\zeta)(g),0)$ by (E31.7); its central coordinate
does not affect conjugation. Conjugation in (E31.2) is therefore
exactly (E31.3).
The character $\kappa_3(\zeta)$ is nonzero: a cube root of
$\zeta$ would be a unit in $\mathbb Z[\zeta]$, whose unit group
$\mu_6$ has only $\pm1$ as cubes.
This proves the asserted nonzero arithmetic cross term.

Now pass to the prime $(\pi)$ over $7$.
Its residue field is $\mathbb F_7$ and
$\zeta\equiv-3\equiv4\pmod\pi$, so
$K_{(\pi)}=\mathbb Q_7$.
The character $\kappa_3(\zeta)$ is unramified and nonzero:
choose $\zeta_9^3=\zeta$ in the unramified cubic extension.
Arithmetic Frobenius sends $\zeta_9$ to $\zeta_9^7$, whose
ratio to $\zeta_9$ is $\zeta^2$.
The Kummer character of the uniformizer $7$ has tame-inertia
value $\zeta$ after a choice of generator, and Frobenius can
be chosen to fix its selected roots.
Consequently (E31.8) sends Frobenius and tame inertia to
$(2,0)$ and $(0,1)$ in $H_3^{\mathrm{ab}}$.
Every pair of lifts has commutator $(0,0,2)\ne1$.

An arithmetic section lifting $\alpha$ would yield a
homomorphism $G_{\mathbb Q_7}\to H_3$, because the standard
target action is trivial. Wild inertia is pro-$7$ and maps
trivially to this finite $3$-group. The tame relation is
$\sigma\tau\sigma^{-1}=\tau^7$.
Every element of $H_3$ has exponent three, so this relation
forces its two images to commute, a contradiction.
The exact field identities, the local leading constants,
all $27$ group elements and all nine choices of the two
central lifts are checked by
\texttt{e\_intrinsic\_obstruction.py}.

The map $\phi_h$ kills $\gamma_3N_3$, $N_3''$ and
$\gamma_3I_3$. Thus it factors through each intrinsic quotient
named in the statement. These factorizations also apply to
the arithmetic maps after restriction to $G_K$, which suffices
to rule out a section over $\mathbb Q$.
There is also a single map over $\mathbb Q$ displaying both
conjugate comparisons: the adjoint of (E31.4) is
\[
 U_3\longrightarrow
 \operatorname{Res}_{K/\mathbb Q}
      (\mathbb P^1\setminus\{0,1,\infty\}).
 \tag{E31.9}
\]
After base change to $\overline{\mathbb Q}$ it is
$(h,\bar h)$, giving the finite geometric image inside
$H_3\times H_3$. The arithmetic map to this image is defined
over $\mathbb Q$ and its geometric kernel contains
$\gamma_3N_3$. Its restriction has the first projection
used above. This supplies an actual global finite quotient
comparison as well.

Adding any element of $3\mathcal K$ to $\alpha$ leaves
(E31.8) unchanged, since the target abelian coefficient group
has exponent three. The same obstruction applies to the
whole coset $\alpha+3\mathcal K$. The zero class is the original
rational tangential section and lifts to the full arithmetic
fundamental group. Both lie over the zero pro-$3$ toric class,
as asserted. In particular this proof distinguishes
liftability inside one specified fiber, using a quotient that
detects an intrinsic direction beyond the inherited toric
comparison of E.28.
\end{proof}
